% Pillar-5 (MIN-REPORT-RL) standalone introduction
\section{Introduction: The Sandwich Problem}
\label{sec:p5-intro}
When a paper reports that ``DAPO outperforms GRPO,'' the named algorithm is a
thin layer sandwiched between two thick, largely undocumented slices. Beneath the
label sits the serving and training substrate: the sampling engine and its
numerics~\citep{kwon2023vllm, zheng2024sglang}, attention kernels and
precision~\citep{dao2022flashattention}, the trainer's tokenization and masking
conventions, and the exact loss implementation shipped by the framework
\citep{vonwerra2022trl, sheng2024verl, hu2024openrlhf, thinkingmachines2024tinker}.
Above the label sits the reward and evaluation apparatus: the verifier and its
answer parser, the held-out split, the prompt template, and whatever
decontamination was (or was not) performed \citep{zhang2024gsm1k,
riddell2024contamination}. The label names the filling; the experiment is the
sandwich. We call the resulting failure of reference the \emph{sandwich problem}:
two papers using the same algorithm label routinely describe different
experiments, and two papers using different labels routinely differ more in their
bread than in their filling.

This is a position paper. Its position is: \textbf{RL-for-LLM results are
stack-conditioned, and venues should require authors to report the stack, not
merely the label.} We do not claim that algorithmic labels are meaningless---the
loss-form differences among GRPO~\citep{shao2024deepseekmath, guo2025deepseekr1},
DAPO~\citep{yu2025dapo}, Dr.GRPO~\citep{tong2025drgrpo}, and
GSPO~\citep{qwen2025gspo} are real and sometimes consequential. We claim that on
the evidence available to us the stack's contribution to reported outcomes is at
least as large as the label's, and that the community's reporting conventions are
calibrated to the wrong term. A $17\times$ outcome swing from swapping only the
backend (\secref{sec:p5-evidence}) dwarfs every label-vs-label gap in our
corpus; a label can even invert its own telemetry when re-implemented on a
different stack. Deep RL went through this reckoning once
already~\citep{henderson2018deep, colas2019hitchhiker, agarwal2021deep};
RL-for-LLMs has re-created the conditions with a larger and less inspectable
stack.

Our contributions are deliberately those of a position paper: evidence, a
standard, a threat model, and enforcement machinery.
\begin{enumerate}[leftmargin=1.5em, itemsep=1pt, topsep=2pt]
\item \textbf{Evidence} (\secref{sec:p5-evidence}): four exhibits---a
$17\times$ backend flip under otherwise identical configuration; a cross-stack
label flip in which ``DAPO'' yields mean ZVF $0.00$ on one stack and $0.58$ on
another; a 12-cell head-to-head in which four labeled algorithms on one fixed
stack land within noise; and cross-library ZVF variation from the companion
Pillar-2 paper of \ours{}.
\item \textbf{The eight-item minimum reportable stack} (\secref{sec:p5-stack}):
a checklist small enough to enforce, in which every item is justified by a
documented lever that can move or flip a reported ranking.
\item \textbf{An RL-reproducibility threat model} (\secref{sec:p5-threat}): a
taxonomy of confound vectors and a flip-risk grading R0--R5 that classifies how
dangerous a given cross-paper comparison is.
\item \textbf{An enforceability toolchain} (\secref{sec:p5-toolchain}): a
specification for a manifest-emitting trainer plugin
(\texttt{trl-min-report-rl}), a manifest differ with CI verdicts
(\texttt{grpo-stackdiff}), a machine-readable stack registry, a 0--100 auditor
badge, and a citation-lever tracer (LeverTrace).
\end{enumerate}
Throughout, we scope claims to their evidence: benchmark results come from
\ours{} (\secref{sec:benchmark}); several exhibits come from internal program
records that we describe as such; and toy-scale theory validations are marked
directional. The paper ends with limitations and a concrete call to action for
venues (\secref{sec:p5-conclusion}).
